% =====================================================================
%  Thème 3 — EXERCICES — Niveau DIFFICILE
% =====================================================================
\documentclass[11pt,a4paper]{article}
\usepackage[utf8]{inputenc}
\usepackage[T1]{fontenc}
\usepackage{lmodern}
\usepackage[french]{babel}
\usepackage{amsmath,amssymb,mathtools}
\usepackage{icomma}
\usepackage{xcolor}
\usepackage[most]{tcolorbox}
\usepackage{enumitem}
\usepackage{array}
\usepackage{booktabs}
\usepackage{fancyhdr}
\usepackage[hidelinks]{hyperref}
\usepackage[a4paper,margin=2.2cm,headheight=26pt,headsep=10pt]{geometry}

\definecolor{primary}{RGB}{223,87,99}
\definecolor{primarydark}{RGB}{150,42,55}

\tcbset{boxbase/.style={enhanced,breakable,boxrule=0.9pt,arc=2.5pt,
   left=3mm,right=3mm,top=2.3mm,bottom=2.3mm,
   fonttitle=\bfseries,coltitle=white,toptitle=0.6mm,bottomtitle=0.6mm}}
\newtcolorbox[auto counter]{exercice}[1][]{boxbase,colback=primary!4,colframe=primary,
  title={\textsc{Exercice}~\thetcbcounter\ifx\relax#1\relax\relax\else\textmd{ --- \itshape #1}\fi}}

\pagestyle{fancy}
\fancyhf{}
\fancyhead[L]{\small\textcolor{primarydark}{\textbf{Fiche 6} — Les limites}}
\fancyhead[R]{\small\textcolor{primarydark}{\textsc{Thème 3 — Difficile}}}
\fancyfoot[C]{\small\thepage}
\renewcommand{\headrulewidth}{0.8pt}
\renewcommand{\headrule}{\hbox to\headwidth{\color{primary}\leaders\hrule height \headrulewidth\hfill}}

\newcommand{\R}{\mathbb{R}}

\begin{document}

\begin{center}
{\Large\bfseries\color{primarydark}Thème 3 — Limites à l'infini : le terme dominant}\\[2pt]
{\large\color{primary}Exercices — Niveau Difficile}
\end{center}
\vspace{4pt}
{\small\itshape Niveau concours : un paramètre à ajuster, une combinaison à réduire, et une mise en
situation avec asymptote horizontale.}
\vspace{6pt}

\begin{exercice}[une famille de fonctions rationnelles]
Pour un réel $m$, on pose $f_m(x)=\dfrac{m x^2+2x-1}{x^2+3}$.
\begin{enumerate}[label=\alph*),leftmargin=2em,itemsep=3pt]
  \item Pour $m=4$, calculer $\displaystyle\lim_{x\to+\infty} f_4(x)$ et donner l'asymptote horizontale.
  \item Exprimer $\displaystyle\lim_{x\to+\infty} f_m(x)$ en fonction de $m$ (on suppose $m\neq 0$).
  \item Déterminer $m$ pour que la courbe admette la droite $y=2$ comme asymptote horizontale.
  \item Existe-t-il une valeur de $m$ pour laquelle $\displaystyle\lim_{x\to+\infty} f_m(x)=0$ ?
        \emph{Indication : que devient le degré du numérateur si $m=0$ ?}
\end{enumerate}
\end{exercice}

\begin{exercice}[réduire puis passer à la limite]
Soit $h(x)=\dfrac{x^2+1}{x-1}-x$, définie sur $\R\setminus\{1\}$.
\begin{enumerate}[label=\alph*),leftmargin=2em,itemsep=3pt]
  \item Réduire $h(x)$ au même dénominateur.
  \item Développer et simplifier le numérateur pour écrire $h(x)$ comme une seule fraction rationnelle.
  \item En déduire $\displaystyle\lim_{x\to+\infty} h(x)$ et $\displaystyle\lim_{x\to-\infty} h(x)$.
  \item La courbe de $h$ admet-elle une asymptote horizontale ? Laquelle ?
\end{enumerate}
\end{exercice}

\begin{exercice}[mise en situation : capacité d'accueil d'un milieu]
Une culture de bactéries compte $N(t)=\dfrac{500\,t}{t+20}$ milliers d'individus, où $t\geq 0$ est le
temps (en heures).
\begin{enumerate}[label=\alph*),leftmargin=2em,itemsep=3pt]
  \item Calculer $N(0)$, $N(20)$ et $N(180)$ (arrondir au millier).
  \item Calculer $\displaystyle\lim_{t\to+\infty} N(t)$ par la règle du terme dominant.
  \item Interpréter cette limite : que représente-t-elle pour la population ? Comment s'appelle la droite
        horizontale correspondante ?
  \item La population dépassera-t-elle un jour $500\,000$ individus ? Justifier à partir de la limite.
\end{enumerate}
\end{exercice}

\end{document}
